%% LyX 2.4.4 created this file.  For more info, see https://www.lyx.org/.
%% Do not edit unless you really know what you are doing.
\documentclass[english]{article}
\usepackage[T1]{fontenc}
\usepackage[latin9]{inputenc}
\usepackage{enumitem}
\usepackage{amsmath}
\usepackage{amsthm}
\usepackage{amssymb}
\usepackage{cancel}
\usepackage{geometry}
\geometry{verbose,tmargin=1in,bmargin=1in,lmargin=1in,rmargin=1in}

\makeatletter
%%%%%%%%%%%%%%%%%%%%%%%%%%%%%% Textclass specific LaTeX commands.
\numberwithin{equation}{section}
\numberwithin{figure}{section}
\newlength{\lyxlabelwidth}      % auxiliary length 

%%%%%%%%%%%%%%%%%%%%%%%%%%%%%% User specified LaTeX commands.
\usepackage{fancyhdr}
\usepackage{lastpage}
\pagestyle{fancy}
\usepackage{enumitem}
\usepackage{ifthen}
\everymath=\expandafter{\the\everymath\displaystyle}
%\DeclareMathSizes{10}{10}{10}{10}
\usepackage{multicol}

\usepackage{tikz}
\usetikzlibrary{patterns}
\usetikzlibrary{plotmarks}
\usepackage{pgfplots}
\pgfplotsset{compat=newest}

\usepackage{advdate}    % Advancing/saving dates
\usepackage{datetime}   % Dates formatting
\usepackage{datenumber} % Counters for dates
\usdate
% Added by lyx2lyx
\setlength{\parskip}{\smallskipamount}
\setlength{\parindent}{0pt}

\makeatother

\usepackage{babel}
\begin{document}
\renewcommand{\author}{Dr.\,James Ashe}

\newcommand{\classNum}{335}

\newcommand{\classSec}{A}

\newcommand{\examType}{Midterm Exam}

\renewcommand{\date}{\formatdate{5}{4}{2013}}

%%First page's header%%%%%%%%%%%%%%%%%%%%%%
\fancypagestyle{firststyle} %%header settings for first pagr
{
	\fancyhead[L]{MTH \classNum{}(\classSec{}) \author{}\\
	\textbf{\examType{}} \date{}}
	\fancyhead[C]{\textbf{Solutions}}
	\setlength{\headsep}{14pt}
}
%%%%%%%%%%%%%%%%%%%%%%%%%%%%%%%%%%%%%%%%%%

%%Next page's header%%%%%%%%%%%%%%%%%%%%%%%
\setlist{nolistsep}
\lhead{\classNum{}(\classSec{})} 
\chead{\textbf{}} 
\cfoot{\thepage{}~of~\pageref{LastPage}} 
\setlength{\headsep}{5pt} 
%%%%%%%%%%%%%%%%%%%%%%%%%%%%%%%%%%%%%%%%%%%

%%Various settings; see the preamble for other settings%%%%%
%\setenumerate[0]{leftmargin=12pt,itemindent=0pt,labelwidth=0pt}
\setlist{topsep=.5pt}
\renewcommand{\labelenumi}{\textbf{\arabic{enumi}.}}
\renewcommand{\labelenumii}{\textbf{\alph{enumii}.}}
\thispagestyle{firststyle}
%%%%%%%%%%%%%%%%%%%%%%%%%%%%%%%%%%%%%%%%%%

\newcommand{\Dspace}{.8cm}\small
\begin{enumerate}[resume]
\item Determine the following:

\begin{enumerate}
\item Find an integer $x$ such that $0\leq x<4$ and $x=15\bmod4$

\begin{enumerate}
\item []$x=3$. Using, $15=3\cdot4+\mathbf{3}$ or $4|15-3$.
\end{enumerate}
\item $\mbox{gcd}(2\cdot3^{40}\cdot5^{20}\cdot13,2\cdot3\cdot5^{19}\cdot7)=2\cdot3\cdot5^{19}$
\item Write the set $A=\left\{ \dots-3,0,3,6,\dots\right\} $ in set-builder
form.

\begin{enumerate}
\item []$A=\left\{ x|x=3n\mbox{ for }n\in\mathbb{Z}\right\} $\vfill
\end{enumerate}
\end{enumerate}
\item Complete the following:

\begin{enumerate}
\item $a\cdot s+b\cdot t=1$ for integers $a,b,s,t\in\mathbb{Z}$. What
is $\gcd(a,b)$?

\begin{enumerate}
\item []$\gcd(a,b)=1$ since $a\cdot s+b\cdot t=1$
\end{enumerate}
\item If $a\cdot s+b\cdot t=1$ is there an integer $d>1$ such that $d|a$
and $d|b$? Explain.

\begin{enumerate}
\item []No. The greatest common divisor is $1$ since $\gcd(a,b)=1$, and
$d>1$.\vfill
\end{enumerate}
\end{enumerate}
\item Complete the following:

\begin{enumerate}
\item Find all the subsets of $A=\left\{ a,b,c\right\} $.

\begin{enumerate}
\item []$\emptyset,\left\{ a\right\} ,\left\{ b\right\} ,\left\{ c\right\} ,\left\{ a,b\right\} ,\left\{ a,c\right\} ,\left\{ b,c\right\} ,\left\{ a,b,c\right\} $
\end{enumerate}
\item If $\left|A\right|=4$ how many subsets does $A$ have?

\begin{enumerate}
\item []$2^{4}=16$\vfill
\end{enumerate}
\end{enumerate}
\item Let $A=\left\{ 1,2\right\} $ and $B=\left\{ 2,4\right\} $.

\begin{enumerate}
\item Find $A\cup B=\left\{ 1,2,4\right\} $
\item Find $A-B=\left\{ 1\right\} $
\item Find $A\times B=\left\{ \left(1,2\right),\left(1,4\right),\left(2,2\right),\left(2,4\right)\right\} $\vfill
\end{enumerate}
\item Prove that $A=A\cup A$.

\begin{proof}
$\left(A\subseteq A\cup A\right)$: Let $x\in A$. $x\in A$ or $x\in A$
so then $x\in A\cup A$. Thus $A\subseteq A\cup A$.

$\left(\left(A\cup A\right)\subseteq A\right)$: Let $x\in A\cup A$.
So $x\in A$ or $x\in A$. Thus $A\cup A\subseteq A$.\vfill
\end{proof}
\item Let $A\subseteq B$ and $B\subseteq C$. Prove that $A\subseteq C$.

\begin{proof}
Let $a\in A$ then , and $a\in B$ implies that $a\in C$. So then
$A\subseteq C$.\vfill
\end{proof}
\item Give the definition of the following:

\begin{enumerate}
\item A \emph{1-1 function}.

$f:A\rightarrow B$ is a function is 1-1 if and only if $f(x_{1})=f(x_{2})$
then $x_{1}=x_{2}$.
\item A \emph{transitive relation.}

A relation $\mathcal{R}$ is \emph{transitive }if $x\mathcal{R}y$
and $yRz$ implies that $yRz$.\vfill
\end{enumerate}
\item Let $f:A\rightarrow B$ be a function from set $A$ to set $B$.

\begin{enumerate}
\item If $\left|A\right|=20$ and $\left|B\right|=21$ can $f$ be onto?

No, since $\left|A\right|<\left|B\right|$. 
\item If $f$ is a constant function and $\left|B\right|>1$, can it be
$1-1$? ($f$ is \emph{constant function }if $f(a)=b$ for all $a\in A$
and some $b\in B$.)

Yes. Let $A=\left\{ a\right\} $ and $f(a)=b$. \emph{In general,
}a constant function is not 1-1.\emph{ }\vfill
\end{enumerate}
\item Let $\mathcal{R}=\left\{ \left(a,b\right),\left(b,a\right),\left(b,c\right),\left(a,a\right),\left(b,b\right),\left(c,c\right)\right\} $
be a relation on the set $E=\left\{ a,b,c\right\} $.

\begin{enumerate}[resume]
\item Is $\mathcal{R}$ reflexive? Explain. 

Yes. $\left(a,a\right),\left(b,b\right),\left(c,c\right)\in\mathcal{R}$.
\item Is $\mathcal{R}$ symmetric? Explain.

No. $\left(b,c\right)\in\mathcal{R}$ but $\left(c,b\right)\notin\mathcal{R}$.\vfill
\end{enumerate}
\item Define the relation $\mathcal{R}$ on the set of even integers as
$x\mathcal{R}y\iff2|x\cdot y$ (i.e., $x\cdot y$ is an even integer).
Show that $\mathcal{R}$ is an equivalence relation.

\begin{proof}
Let $x,y,z\in2\mathbb{Z}$, i,e., $x=2n_{0},y=2n_{1},$ and $z=2n_{2}$
for some $n_{0},n_{1},n_{2}\in\mathbb{Z}$.\\
Reflexive: $x\cdot x=2n\cdot2n=4n^{2}$ and $2|4n^{2}$ so $x\mathcal{R}x$.

Symmetric: Suppose that $2|x\cdot y$. $x\cdot y=y\cdot x$ so then
$2|yx$.

Transitive: Suppose that $2|x\cdot y$ and $2|y\cdot z$. $x\cdot z=2n_{0}2n_{2}=4n_{0}n_{2}$
so then $2|x\cdot z$.

$\mathcal{R}$ is reflexive, symmetric, and transitive so it is an
equivalence relation.
\end{proof}
\vfill
\end{enumerate}

\end{document}
